%! Equivalent Representations of Polynomial Expressions — Unit 2, Lesson 2.6 — Exit Ticket
\documentclass[10pt]{article}
\usepackage{precalculus-article}
\usepackage{precalculus-boxes}
\newcommand{\TallMath}[1]{$\displaystyle #1\rule[-1.4em]{0pt}{3.2em}$}

\begin{document}

\pageheader{Unit 2, Lesson 2.6}{Exit Ticket}

\vspace{0.12in}

\begin{enumerate}[itemsep=12pt]

\item These are two forms of the same function:
      $q(x) = (x+1)(x-5) = x^2 - 4x - 5$. For each question, name the form you
      would use, then answer it.

\begin{enumerate}[label=(\alph*), itemsep=6pt]
  \item $x$-intercepts. \quad Form: \blank{2.4cm} \quad
        Answer: \blank{3.2cm}
  \item $\lim_{x \to \infty} q(x)$. \quad Form: \blank{2.4cm} \quad
        Answer: \blank{2cm}
  \item $y$-intercept. \quad Form: \blank{2.4cm} \quad
        Answer: \blank{2cm}
\end{enumerate}

\item Use row 3 of Pascal's Triangle ($1, 3, 3, 1$) to expand $(x+2)^3$.

\begin{work}
  (x+2)^3 &= x^3 + 3x^2(2) + 3x(2)^2 + (2)^3 \\
          &= x^3 + 6x^2 + 12x + 8
\end{work}

$(x+2)^3 = $ \blank{6cm}

\item Divide $x^2 + 5x + 1$ by $x + 2$.

\begin{work}
  &\begin{array}{r@{\;}rrr}
              &    x  &  +3  &     \\ \cline{2-4}
   x+2\,\big) &   x^2 & +5x  & +1  \\
   -          &   x^2 & +2x  &     \\ \cline{2-4}
              &       &  3x  & +1  \\
   -          &       &  3x  & +6  \\ \cline{2-4}
              &       &      & -5
  \end{array}
\end{work}

\begin{enumerate}[label=(\alph*), itemsep=6pt]
  \item Quotient: \blank{2.6cm} \qquad Remainder: \blank{1.6cm}
  \item Is $x+2$ a factor of $x^2 + 5x + 1$? \quad \blank{1.6cm} \quad
        How do you know? \blank{5cm}
\end{enumerate}

\end{enumerate}

\end{document}
